% Generated from locale/tr/content/proof-theory/cut-elimination/ce-largest.tex; do not hand-edit.


\olfileid[tr]{pt}{cut}{inv}

\olsection{En Büyük Kesmelerin Kaldırılması}

\olref[seq][inv]{lem:G3c-invert}'in ispatı, \Log{G3c}'nin bir mantıksal
kuralın (\RightR{\lexists} ile \LeftR{\lforall} dışında) sonucunu
!!{prove}{}dığında, kuralın öncüllerinin her birini de !!{prove}{}dığını ve
bunu ispatın !!{height}~değerini artırmadan yaptığını göstermişti. Daha iyisini
yapabilir ve bunun
$\Log{G3c} + \CutCS$ için de geçerli olduğunu gösterebiliriz.

Daha önce olduğu gibi, bir \CutCS{} çıkarımının \emph{kesme derecesini},
kesme !!{formula}{}ünün derinliği olarak alırız.

\begin{lem}\ollabel{lem:inv-G3c-cut}
$\Log{G3c} + \CutCS$, \RightR{\lexists} ya da \LeftR{\lforall} dışındaki
bir mantıksal kural~$R$'nin sonucunu !!{prove}{}makta ve bunun için
!!a{proof}~$\pi$ kullanıyorsa, her öncülü de !!{prove}{}r: bunu
!!{proof}~$\pi'$ ile (ya da kuralın bir veya iki öncülü olmasına göre
!!{proof}s~$\pi'$, $\pi''$ ile) yapar.
Ayrıca (a)~$\pheight{\pi'}\le\pheight{\pi}$ (ve
$\pheight{\pi''}\le\pheight{\pi}$) ve (b)~\CutCS{} çıkarımlarının sayısı
ile dereceleri,~$\pi'$ içinde (ve $\pi''$ içinde)~$\pi$ içindekilerle
aynıdır.
\end{lem}

\begin{proof}
\cref{pt:seq:inv:lem:G3c-invert,pt:seq:inv:lem:invert-quant}
ispatlarının incelenmesiyle. O ispat,~$\pheight{\pi}$ üzerine tümevarımla
yapılmıştı. Taban durumunda bir aksiyomu başka bir aksiyomla değiştirdik.
Dolayısıyla (a) ve~(b) koşulları önemsiz biçimde sağlanır. Eğer~$\pi$'nin
son kuralı~$R$ ise, onun dolaysız alt-!!{proof}s{}ı istediğimiz
!!{proof}s~$\pi_i$'dir ve hem (a) hem~(b) sağlanır. Öteki bütün durumlarda,
tümevarım varsayımını~$\pi$'nin dolaysız alt-!!{proof}s{}ına, yani son
çıkarım~$R$'nin öncülüne (öncüllerine) götürenlere uyguladık ve ardından
sonuçlara aynı $R$ kuralını uyguladık. Yeni ispatın dolaysız
alt-!!{proof}s{}ı için (a) ve~(b) koşulları geçerlidir; dolayısıyla tüm ispat
için de (a) ve~(b) geçerlidir. Son çıkarımın~\CutCS{} olduğu durumu
doğrulamak kalır. Bunu yine yalnızca tek bir örnek için yapıyoruz: $\pi$'nin
\LeftR{\land}'ın $!B \land !C, \Gamma \Sequent \Delta$ sonucunun
!!a{proof}{}ı olduğunu ve~\CutCS{} ile bittiğini varsayalım:
\[
\AxiomC{}
\RightLabel{$\pi_1$}
\Deduce$!B \land !C, \Gamma \fCenter \Delta, !A$
\AxiomC{}
\RightLabel{$\pi_2$}
\Deduce$!A, !B \land !C, \Gamma \fCenter \Delta$
\RightLabel{\CutCS}
\BinaryInf$!B \land !C, \Gamma \fCenter \Delta$
\DisplayProof
\]
Tümevarım varsayımı, (aynı zamanda $\LeftR{\land}$ sonuç örneklerinin
!!{proof}s{}ı olan) $\pi_1$ ile~$\pi_2$'ye uygulanır ve sırasıyla $!B, !C,
\Gamma \fCenter \Delta, !A$ ile $!A, !B, !C, \Gamma \fCenter \Delta$'nın
!!{proof}s{}ı olan $\pi_1'$ ile~$\pi_2'$'yi verir. Bunları \CutCS{} kullanarak
birleştirip~$\pi'$'yü elde ederiz:
\[
\AxiomC{}
\RightLabel{$\pi_1'$}
\Deduce$!B, !C, \Gamma \fCenter \Delta, !A$
\AxiomC{}
\RightLabel{$\pi_2'$}
\Deduce$!A, !B, !C, \Gamma \fCenter \Delta$
\RightLabel{\CutCS}
\BinaryInf$!B, !C, \Gamma \fCenter \Delta$
\DisplayProof
\]
Açıkça (a) ve~(b) sağlanır.
\end{proof}

\CutCS{} yerine \Cut{} kullansaydık bunun işlemeyeceğine dikkat edin;
çünkü o zaman son !!{proof}{}ın sonucu, öncülde iki $!B, !C, \Gamma$ kopyası
ve ardılda iki $\Delta$ kopyası içerirdi. Büzülmenin kabul edilebilirliğine
başvurmamız gerekirdi; ancak \olref[seq][inv]{prop:G3c-cont-adm}'in ispatı
tersinirlik lemmasını kullanır. Böylece bir döngü içinde akıl yürütmüş
olurduk.

Kesme-giderimine alternatif bir ispat vermek için
\cref{pt:cut:inv:lem:inv-G3c-cut}'i kullanabiliriz. Bu ispatta \CutCS{}
çıkarımlarının en üst oluşumlarını art arda kaldırmak yerine, maksimal
dereceli \CutCS{} çıkarımlarının oluşumlarını kaldırırız. Yani $\Log{G3c} +
\CutCS$ içinde !!a{proof}~$\pi$ ile başlar ve onun herhangi bir yerindeki
bir kesmeyi kaldırırız (onu daha küçük dereceli kesmelerle değiştirebiliriz)
---sonra hiç kesme kalmayana dek bunu sürdürürüz. Tek koşul, ele aldığımız
kesmenin üstünde aynı veya daha yüksek dereceli hiçbir kesmenin
bulunmamasıdır.

\begin{lem}\ollabel{lem:max-cut-red-G3c} $\pi$, $\Log{G3c}+\CutCS$ içinde,
derecesi~$n$ olan bir~\CutCS{} çıkarımıyla biten ve bunun dışında yalnızca
~$\Log{G3c}$ kurallarını ve derecesi $<n$ olan \CutCS{} çıkarımlarını
kullanan !!a{proof} ise, aynı son-sekantın~$\Log{G3c} + \CutCS$ içinde bütün
\CutCS{} çıkarımlarının derecesinin~$<n$ olduğu bir ispatı~$\pi'$ vardır.
\end{lem}

\begin{proof}
!!{proof}~$\pi$ şöyle biter:
\[
\AxiomC{}
\RightLabel{$\pi_1$}
\Deduce$\Gamma \fCenter \Delta, !A$
\AxiomC{}
\RightLabel{$\pi_2$}
\Deduce$!A, \Gamma \fCenter \Delta$
\RightLabel{\CutCS}
\BinaryInf$\Gamma \fCenter \Delta$
\DisplayProof
\]
~$!A$'nın biçimine göre durumları ayırırız.

$!A$'nın atomik olduğunu varsayalım. $\pi_1$ ile $\pi_2$ içindeki bütün
\CutCS{} çıkarımlarının derecesinin $< \depth{!A} = 0$ olması koşulu
gereğince, $\pi_1$ ya da~$\pi_2$ içinde hiçbir \CutCS{} çıkarımı
bulunamayacağına dikkat edin. $\Gamma \Sequent \Delta$'nın kesmesiz bir
ispatı olduğunu $\pheight{\pi_2}$ üzerine tümevarımla gösteririz.
$\pheight{\pi_2} = 0$ ise $!A, \Gamma \Sequent \Delta$ bir aksiyomdur.
$!A$ esas değilse, hem~$\Gamma$'nın hem $\Delta$'nın !!a{element}{}ı olan
atomik bir $!C$ vardır ve $\Gamma \Sequent \Delta$'nın kendisi bir
aksiyomdur. $!A$ esas ise $\Delta = \Delta', !A$'dır. Bu durumda $\pi_1$,
$\Gamma \Sequent \Delta', !A, !A$ ile biter. Büzülmenin kabul
edilebilirliğini veren \olref[seq][inv]{prop:G3c-cont-adm}'i uygulayarak
$\Gamma \Sequent \Delta', !A$'nın kesmesiz bir ispatını elde ederiz.
$\pheight{\pi_2} > 0$ ise bir $R$ kuralıyla biter. Bu kuralın bir öncülünü
alın: öncül $!A, \Gamma', \Pi \Sequent \Delta', \Lambda$ biçimindedir;
burada $\Pi$ ile $\Lambda$, $R$'nin etkin !!{formula}s{}idir.
~$\pi_1$'e ve $\pi_2$'nin ilgili dolaysız alt-ispatına uygulanan
\olref[seq][adm]{prop:weak-G3c-adm}, $\Gamma, \Gamma', \Pi \Sequent
\Delta, \Delta', \Lambda, !A$ ile $!A, \Gamma, \Gamma', \Pi \Sequent
\Delta, \Delta', \Lambda$'nın !!{proof}s{}ını verir ve tümevarım varsayımı
bunlara uygulanır. Böylece $R$'nin her öncülü için $\Gamma, \Gamma', \Pi
\Sequent \Delta, \Delta', \Lambda$'nın !!a{proof}{}ını elde ederiz. $R$'yi
uygulamak $\Gamma, \Gamma \Sequent \Delta, \Delta$'yı verir ve
\olref[seq][inv]{prop:G3c-cont-adm}, $\Gamma \Sequent \Delta$'nın kesmesiz
ispatını sağlar.

$!A$ atomik değilse, !!{main
operator} bakımından durumları ayırırız.
Her durumda \olref{lem:inv-G3c-cut}'i uygulayarak, esas formülü $!A$ olan
sol ve sağ kuralların öncüllerinin !!{proof}s{}ını elde ederiz. Ardından
\olref[top]{lem:cut-adm-G3c} ispatının E durumundaki gibi ilerleriz.
\begin{enumerate}
  \item $!A \ident \lnot !B$. !!{proof}s $\pi_1$ ile $\pi_2$, $\Gamma
  \Sequent \Delta, \lnot !B$ ve $\lnot !B, \Gamma \Sequent \Delta$ ile
  biter. \olref{lem:inv-G3c-cut} gereğince $!B, \Gamma \Sequent \Delta$
  ile $\Gamma \Sequent \Delta, !B$'nin !!{proof}s{}ı olan $\pi_1'$ ile
  $\pi_2'$ vardır. !!{proof}~$\pi'$ şudur:
  \[
  \AxiomC{}
  \RightLabel{$\pi_2'$}
  \Deduce$\Gamma \fCenter \Delta, !B$
  \AxiomC{}
  \RightLabel{$\pi_1'$}
  \Deduce$!B, \Gamma \fCenter \Delta$
  \RightLabel{\CutCS}
  \BinaryInf$\Gamma \fCenter \Delta$
  \DisplayProof
  \]
  \item $!A \ident (!B \lor !C)$. !!{proof}s $\pi_1$ ile $\pi_2$,
  $\Gamma \Sequent \Delta, !B \lor !C$ ve $!B \lor !C, \Gamma \Sequent
  \Delta$ ile biter. \olref{lem:inv-G3c-cut} gereğince (~$\pi_1$'e
  uygulanan \RightR{\lor} tersinimi), $\Gamma \Sequent \Delta, !B,
  !C$'nin !!a{proof}{}ı $\pi_1'$ vardır. \olref{lem:inv-G3c-cut} gereğince
  (~$\pi_2$'ye uygulanan \LeftR{\lor} tersinimi), $!B, \Gamma \Sequent
  \Delta$'nın !!a{proof}{}ı $\pi_2'$ ile $!C, \Gamma \Sequent \Delta$'nın
  !!a{proof}{}ı $\pi_2''$ vardır. $\pi_2''$'ye
  \olref[seq][adm]{prop:weak-G3c-adm}'i uygulayarak $!C, \Gamma \Sequent
  \Delta, !B$'nin !!a{proof}{}ı~$\pi_2'''$'yü de elde ederiz.
  !!{proof}~$\pi'$ şudur:
  \[
  \AxiomC{}
  \RightLabel{$\pi_1'$}
  \Deduce$\Gamma \fCenter \Delta, !B, !C$
  \AxiomC{}
  \RightLabel{$\pi_2'''$}
  \Deduce$!C, \Gamma \fCenter \Delta, !B$
  \RightLabel{\CutCS}
  \BinaryInf$\Gamma \fCenter \Delta, !B$
  \AxiomC{}
  \RightLabel{$\pi_2'$}
  \Deduce$!B, \Gamma \fCenter \Delta$
  \RightLabel{\CutCS}
  \BinaryInf$\Gamma \fCenter \Delta$
  \DisplayProof
  \]
  \item $!A \ident (!B \land !C)$. Alıştırma.
  \item $!A \ident (!B \lif !C)$. !!{proof}s $\pi_1$ ile $\pi_2$,
  $\Gamma \Sequent \Delta, !B \lif !C$ ve $!B \lif !C, \Gamma \Sequent
  \Delta$ ile biter. \olref{lem:inv-G3c-cut} gereğince (~$\pi_1$'e
  uygulanan \RightR{\lif} tersinimi), $!B, \Gamma \Sequent \Delta,
  !C$'nin !!a{proof}{}ı $\pi_1'$ vardır. \olref{lem:inv-G3c-cut} gereğince
  (~$\pi_2$'ye uygulanan \LeftR{\lif} tersinimi), $\Gamma \Sequent
  \Delta, !B$'nin !!a{proof}{}ı $\pi_2'$ ile $!C, \Gamma \Sequent
  \Delta$'nın !!a{proof}{}ı $\pi_2''$ vardır. $\pi_2''$'ye
  \olref[seq][adm]{prop:weak-G3c-adm}'i uygulayarak $!C, !B, \Gamma \Sequent
  \Delta$'nın !!a{proof}{}ı~$\pi_2'''$'yü de elde ederiz.
  !!{proof}~$\pi'$ şudur:
  \[
  \AxiomC{}
  \RightLabel{$\pi_2'$}
  \Deduce$\Gamma \fCenter \Delta, !B$
  \AxiomC{}
  \RightLabel{$\pi_1'$}
  \Deduce$!B, \Gamma \fCenter \Delta, !C$
  \AxiomC{}
  \RightLabel{$\pi_2'''$}
  \Deduce$!C, !B, \Gamma \fCenter \Delta$
  \RightLabel{\CutCS}
  \BinaryInf$!B, \Gamma \fCenter \Delta$
  \RightLabel{\CutCS}
  \BinaryInf$\Gamma \fCenter \Delta$
  \DisplayProof
  \]
  \item $!A \ident \lforall[x][!B(x)]$. !!{proof}s $\pi_1$ ile
  $\pi_2$, $\Gamma \Sequent \Delta, \lforall[x][!B(x)]$ ve
  $\lforall[x][!B(x)], \Gamma \Sequent \Delta$ ile biter.
  \olref{lem:inv-G3c-cut} gereğince $\Gamma \Sequent \Delta, !B(c)$'nin
  !!a{proof}{}ı~$\pi_1'$ vardır.

  Şimdi !!{proof}~$\pi_2$'yi ele alın. Bu ispat
  $\lforall[x][!B(x)], \Gamma \Sequent \Delta$ ile biter. ~$\pi_2$'nin
  son-sekantından yukarı doğru giderek, bir sekantın solunda bulunan ve
  $\pi_2$'nin son-sekantındaki $\lforall[x][!B(x)]$ oluşumuna ``götüren''
  her $\lforall[x][!B(x)]$ oluşumunu işaretleyin; yani:
  (a)~$\pi_2$'nin son-sekantındaki $\lforall[x][!B(x)]$'i işaretleyin.
  (b)~$\lforall[x][!B(x)]$ bir kuralın sonucunun bağlamında oluşuyor ve
  işaretliyse, kuralın öncülündeki (öncüllerindeki) karşılık gelen oluşumu
  işaretleyin. (c)~$\lforall[x][!B(x)]$, bir \LeftR{\lforall} kuralının
  sonucunda esas olup işaretliyse, öncüldeki karşılık gelen etkin
  $\lforall[x][!B(x)]$ ile $!B(t)$ oluşumlarını işaretleyin.
  
  Şimdi bu işaretli $\lforall[x][!B(x)]$ oluşumlarının tümünü ele alın.
  Bunların hiçbiri \LeftR{\lforall} dışında bir kuralda etkin değildir
  (yalnızca \LeftR{\lforall}'ın etkin !!{formula}{}ü işaretlenir).
  $\pi_2$ içindeki bütün işaretli $\lforall[x][!B(x)]$ oluşumlarını silin.
  Bu doğru bir !!{proof} ise, işaretli $\lforall[x][!B(x)]$ oluşumları hiçbir
  zaman etkin olmamıştır; hepsi doğrudan aksiyomlardan gelmiştir.
  !!{proof},~$\Gamma \Sequent \Delta$ ile biter ve $\pi$'yi onunla
  değiştirebiliriz. Maksimal dereceli bir kesmeyi kaldırmış oluruz.

  Aksi durumda elimizdeki doğru bir ispat değildir; çünkü bazı
  \LeftR{\lforall} çıkarımlarında etkin olan !!{formula}s
  $\lforall[x][!B(x)]$'i silmişizdir. Bunları şu biçimdeki ``çıkarımlara''
  dönüştürmüşüzdür:
  \[
  \AxiomC{}
  \RightLabel{$\pi_t$}
  \Deduce$!B(t), \Pi \fCenter \Lambda$
  \RightLabel{\LeftR{\lforall}}
  \UnaryInf$\Pi \fCenter \Lambda$
  \DisplayProof
  \]
  Bu tür en üst ``çıkarımların'' her birini şununla değiştirin:
  \[
  \AxiomC{}
  \RightLabel{$\Subst{\pi_1'}{t}{c}[\Pi \Sequent \Lambda]$}
  \Deduce$\Gamma, \Pi \fCenter \Delta, \Lambda, !B(t)$
  \AxiomC{}
  \RightLabel{$\pi_t[\Gamma \Sequent \Delta]$}
  \Deduce$!B(t), \Gamma, \Pi \fCenter \Delta, \Lambda$
  \RightLabel{\CutCS}
  \BinaryInf$\Gamma, \Pi \fCenter \Delta, \Lambda$
  \DisplayProof
  \]
  ve onun altındaki her sekantın soluna $\Gamma$'yı, sağına $\Delta$'yı
  ekleyin. Öncülünde işaretli bir $!B(t)$ bulunan bütün \LeftR{\lforall}
  ``çıkarımları'', bir $!B(t)$ üzerindeki \CutCS{} çıkarımıyla değiştirilene
  dek tekrarlayın. Ortaya çıkan !!{proof}{}ın (artık doğru bir !!{proof}{}tır)
  son-sekantı $\Gamma, \dots, \Gamma \Sequent \Delta, \dots, \Delta$
  biçimindedir. Büzülmenin kabul edilebilirliğini uygulayıp bunu $\Gamma
  \Sequent \Delta$'nın !!a{proof}{}ına dönüştürün ve $\pi$'yi onunla
  değiştirin. $\lforall[x][!B(x)]$ üzerindeki bir kesmeyi $!B(t)$ üzerindeki
  (muhtemelen çok sayıda) kesmeyle değiştirmiş oluruz; ancak bunların hepsi
  daha düşük derecelidir. $\pi_1$ ya da $\pi_2$ içinde zaten bulunan bütün
  kesmeler korunur; fakat onların da hepsi daha düşük derecelidir.
\end{enumerate}
\end{proof}

\begin{prob}
\olref{lem:max-cut-red-G3c} ispatında $!A \ident (!B \land !C)$ olan durumu
doğrulayın. Yani yalnızca derecesi $<\depth{!B \land !C}$ olan kesmeler
içeren $\pi_1$ ile $\pi_2$'den oluşan ve
\[
\AxiomC{}
\RightLabel{$\pi_1$}
\Deduce$\Gamma \fCenter \Delta, !B \land !C$
\AxiomC{}
\RightLabel{$\pi_2$}
\Deduce$!B \land !C, \Gamma \fCenter \Delta$
\RightLabel{\CutCS}
\BinaryInf$\Gamma \fCenter \Delta$
\DisplayProof
\]
ile biten !!a{proof}~$\pi$ varsa, yalnızca derecesi $<\depth{!B \land !C}$
olan kesmeler içeren bir $\Gamma \Sequent \Delta$ !!a{proof}{}ı~$\pi'$
bulunduğunu gösterin.
\end{prob}

% \begin{prob}
% \olref{lem:inv-G3c-cut} ispatında $!A \ident \lexists[x][!B(x)]$ olan
% durumu doğrulayın.
% \end{prob}


\olref{lem:max-cut-red-G3c}, !!a{proof} içindeki en yüksek dereceli bir
\CutCS{} çıkarımını daha düşük dereceli \CutCS{} çıkarımlarıyla
değiştirebileceğimizi saptar. Bu lemmayı, herhangi bir sayıdaki \CutCS{}
çıkarımını !!a{proof}{}tan, maksimal~\CutCS{} ile biten en üst
alt-!!{proof}s{}ına art arda uygulayarak kaldırabileceğimizi göstermek için
yine kullanabiliriz.

\begin{thm}\ollabel{thm:cut-elim-G3c-largest}
$\Log{G3c} + \CutCS$ içindeki her !!{proof} $\pi$,~\Log{G3c} içinde bir
ispata dönüştürülebilir.
\end{thm}

\begin{proof}
  Önce~$\pi$ içindeki maksimal dereceli bütün kesmelerin
  kaldırılabileceğini ispatlarız. $d$,~$\pi$ içinde oluşan kesmelerin
  maksimal derecesi ve $n$, derecesi~$d$ olan kesmelerin sayısı olsun.
  İspat~$n$ üzerine tümevarımladır. $n=0$ ise ispatlanacak bir şey yoktur.
  $n>0$ ise derecesi~$d$ olan en üst bir kesme seçin ve onunla biten
  alt-!!{proof}{}ı $\pi_1$ olarak alın. \olref{lem:max-cut-red-G3c} gereğince,
  aynı son-sekantın bütün kesmelerinin derecesi $<d$ olan
  !!a{proof}{}ı~$\pi_1'$ vardır. ~ $\pi$ içindeki $\pi_1$'i $\pi_1'$ ile
  değiştirin. Bunun sonucunda derecesi~$d$ olan $n-1$ kesme içeren
  !!a{proof} elde edilir ve tümevarım varsayımı uygulanır.

  Şimdi sonuç~$d$ üzerine tümevarımla çıkar. $d=0$ ise bütün kesmeler
  atomiktir ve önceki sonuç gereğince hepsi kaldırılabilir. $d>0$ ise
  önceki sonuç, derecesi $d$ olan bütün kesmelerin derecesi~$<d$ olan
  kesmelerle değiştirildiği bir ispat verir. $d$,~$\pi$ içindeki kesmelerin
  maksimal derecesi olduğundan, böylece kesmelerin maksimal derecesinin
  $<d$ olduğu bir ispatımız vardır ve tümevarım varsayımı uygulanır.
\end{proof}

